\documentclass[12pt]{article}

\usepackage{amsmath,amssymb}
\usepackage{fullpage}
\usepackage{enumitem}
\usepackage{hyperref}
\usepackage{listings}
\usepackage{xcolor}

\setlength{\parindent}{0pt}

\lstset{
    language=Python,
    basicstyle=\ttfamily\small,
    keywordstyle=\color{blue},
    commentstyle=\color{gray},
    stringstyle=\color{teal},
    frame=single,
    columns=fullflexible,
    keepspaces=true,
    showstringspaces=false
}

\begin{document}

\begin{center}
	\hrule
	\vspace{0.4cm}
	{\textbf{\large Master Theorem Simplified}}\\[0.2cm]
	INF1008 Data Structures and Algorithms
\end{center}

\textbf{Name:} Timothy Chia\hspace{\fill} \textbf{Student ID:} 2501530\\

\hrule

\vspace{0.3cm}

\paragraph{The Master Theorem.} Divide-and-conquer algorithms are commonly analysed using recurrence relations. A large class of such recurrences can be solved using the \textit{Master Theorem}, which provides asymptotic bounds without requiring full expansion of the recurrence.

\vspace{0.2cm}

\subsection*{Problem Form}

The Master Theorem applies to recurrences of the form:
\[
	T(n) =
	\begin{cases}
		\Theta(1),                                      & n < c_1   \\
		a \, T\!\left(\frac{n}{b}\right) + \Theta(n^i), & n \ge c_1
	\end{cases}
\]

where:
\begin{itemize}[leftmargin=1.2cm]
	\item $a \ge 1$ is the number of subproblems,
	\item $b > 1$ is the factor by which the input size is reduced,
	\item $\Theta(n^i)$ represents the non-recursive work done at each level,
	\item $c_1$ is a constant defining the base case threshold.
\end{itemize}

This recurrence models algorithms that split a problem into smaller subproblems,
solve them recursively, and perform additional work to combine the results.

\vspace{0.2cm}

\subsection*{Recursion Tree Interpretation}

The recurrence can be visualised as a recursion tree:
\begin{itemize}[leftmargin=1.2cm]
	\item The height of the tree is $\log_b n$.
	\item Each level contains $a^k$ subproblems of size $n / b^k$.
	\item The total work at each level depends on both the number of subproblems
	      and the cost per subproblem.
\end{itemize}

The Master Theorem determines where the dominant contribution to the total runtime occurs: at the root, evenly across all levels, or at the leaves.

\vspace{0.2cm}

\subsection*{Asymptotic Cases}

Let $b^i$ represent the growth rate of the non-recursive work per level.

\subsubsection*{Case 1: $a > b^i$ (Leaf-Dominated)}

If the number of subproblems grows faster than the work per level:
\[
	T(n) = \Theta\!\left(n^{\log_b a}\right)
\]

In this case, the total work increases as the recursion proceeds downward,
and the leaves of the recursion tree dominate the overall complexity.

\vspace{0.15cm}

\subsubsection*{Case 2: $a = b^i$ (Balanced Contribution)}

If the recursive and non-recursive work grow at the same rate:
\[
	T(n) = \Theta(n^i \log n)
\]

Here, each level of the recursion tree contributes an equal amount of work.
The total runtime is the product of the work per level and the height of the tree.

\vspace{0.15cm}

\subsubsection*{Case 3: $a < b^i$ (Root-Dominated)}

If the non-recursive work dominates:
\[
	T(n) = \Theta(n^i)
\]

The majority of the work is performed near the root of the recursion tree,
and the cost of deeper recursive calls becomes negligible.

\vspace{0.2cm}

\subsection*{Summary}

This simplified definition provides a fast and reliable method for analysing
recurrences by comparing the growth rate of recursive calls
($a$) against the growth rate of non-recursive work ($b^i$).
By identifying which part of the recursion tree dominates the computation,
the asymptotic time complexity can be determined directly.

\end{document}
